\section{Implementaci\'on}

La implementaci\'on de la librer\'ia se encuentra en el archivo de cabecera
\verb|jumpsuit.h|. Esta librer\'ia es \textit{header-only} 
(o \textit{stb-style} por las librer\'ias de Sean Barret \cite{stbnothings}).

Siguiendo una filosof\'ia de desarrollo de bajo nivel, se ha buscado minimizar
las dependencias externas, utilizando principalmente la librer\'ia est\'andar
de C y \verb|cblas.h| 
(espec\'ificamente OpenBLAS\cite{openblas}, dado que es open-source)
para optimizar operaciones matriciales.

\subsection{Estructuras de Datos Principales}

La librer\'ia gira en torno a la estructura \verb|IndexPQ|,
la cual mantiene el estado del \'indice:

\begin{lstlisting}[language=C, caption=Estructura IndexPQ]
typedef struct {
    int dimension;              // = d
    int n_subvectors;           // = m
    int n_bits_per_value;       // = log_2(k*)

    int n_iter;

    ////////////////////////////////////////////////////
    // Calculated automatically based on init params:
    
    int centroids_per_page;     // k*
    int subvector_dimension;

    ////////////////////////////////////////////////////

    ////////////////////////////////////////////////////
    // Filled on train and add
    
    int n_vectors;              // = n

                                //   (C_i = [c_1, c_2, ..., c_k*])
    float *codebook;            // = [C_1, C_2, ..., C_m]

                                //   (Idx_i = [idx_1, idx_2, ..., idx_m])
    int *quantized_codes;       // = [Idx_1, Idx_2, ..., Idx_n]
                                
    ////////////////////////////////////////////////////
} IndexPQ;
\end{lstlisting}

\newpage

\begin{lstlisting}[language=C, caption=Estructura IndexPQ SearchResult]
typedef struct {
    int n_vectors;
    int n_neighbours;

    int *indices;
    float *distances;
} IndexPQ_SearchResult;
\end{lstlisting}

\newpage

\subsection{API e Implementaci\'on de Funciones}

\subsubsection{Inicializaci\'on}

\begin{lstlisting}[language=C, caption=Inicializaci\'on]
IndexPQ index_pq_init(
  int dimension,
  int n_subvectors,
  int n_bits_per_value)
\end{lstlisting}

Esta funci\'on configura los hiperpar\'ametros del algoritmo.
Calcula autom\'aticamente la dimensi\'on de los subvectores ($D/m$)
y la cantidad de centroides necesarios por subespacio ($2^{bits}$).

Adem\'as, establece por defecto 25 iteraciones para el proceso de entrenamiento
de K-Means, el cual se puede cambiar manualmente antes de comenzarlo.

\subsubsection{Entrenamiento:}

\begin{lstlisting}[language=C, caption=Entrenamiento]
void index_pq_train(IndexPQ *index, float *vectors, int n_vectors)
\end{lstlisting}

Su objetivo es generar el \verb|codebook| (centroides) para cada subespacio.
El proceso sigue estos pasos:
\begin{enumerate}
  \item \textbf{Segregaci\'on}: 
    Los vectores de entrada se copian en arreglos contiguos de subvectores
    para mejorar la localidad de datos.

  \item \textbf{K-Means++}: 
    Para cada subespacio, se inicializan los centroides utilizando el algoritmo 
    K-Means++\cite{wiki:K-means}, lo que mejora la convergencia respecto a una
    inicializaci\'on puramente aleatoria.

  % TODO(fede): Mostrar resultados de profiling
  \item \textbf{Optimizaci\'on BLAS}:
    El c\'alculo de distancias es uno acelerado, donde la 
    distancia euclideana (\(L^2\)) se calcula\cite{samuelalbanie}:

\begin{align}
    d^2(p, q) &= \sum_{i=1}^{n}{(q_i - p_i)^2} 
    = ||p||_2^2 + ||q||_2^2 - 2p^Tq
\end{align}

    permitiendo el calculo prematuro de las normas, y luego el producto 
    matricial de todos los vectores que se tengan que comparar, que en el 
    calculo de los centroides son multiples pasadas de 
    \(\textbf{subvectores}\times\textbf{centroides}\).
    El producto matricial durante el clustering se realiza en 
    \verb|oj__multiply_vectors_blas|, utilizando la funci'on \verb|sgemm| de
    BLAS. 

    Adem\'as, se utiliza \(d^2\), ya que 
    \(d_1 > d_2\) y \(d_1^2 > d_2^2\) son equivalentes para todo \(d_i\) real.
    

\end{enumerate}

\subsubsection{Adici\'on de Vectores}

\begin{lstlisting}[language=C, caption=Agregar Vectores al \'Indice]
void index_pq_add(IndexPQ *index, float *vectors, int n_vectors)
\end{lstlisting}

Una vez entrenado el modelo, esta funci\'on cuantifica los vectores originales.
Para cada subvector, busca el centroide m\'as cercano en su respectivo
subespacio utilizando \verb|oj__get_closest_centroid| y almacena exclusivamente
el \'indice entero de dicho centroide.

\subsubsection{B\'usqueda Vectorial}

\begin{lstlisting}[language=C, caption=Busqueda Vectorial]
IndexPQ_SearchResult index_pq_search(
  IndexPQ *index,
  float *vectors,
  int n_vectors,
  int n_neighbours)
\end{lstlisting}

Implementa el m\'etodo de Distancia Asim\'etrica (ADC).
En lugar de decodificar los vectores almacenados, el algoritmo:
\begin{enumerate}
  \item Genera una \textbf{tabla de b\'usqueda de distancias} 
    entre el vector de consulta y todos los centroides del \textit{codebook}.
  \item Calcula la distancia aproximada sumando los valores pre-calculados
    en la tabla bas\'andose en los \'indices almacenados.
  \item Mantiene un conjunto de los $k$ vecinos m\'as cercanos mediante una
    inserci\'on ordenada.
\end{enumerate}

\subsection{Gesti\'on de Memoria y Portabilidad}

Se utiliz\'o asignaci\'on de memoria alineada (\verb|oj__aligned_alloc|) donde 
fuese necesario para poder ultilizar \verb|sgemm|, y mantener una velocidad de
ejecuci\'on razonable.


